\begin{figure} [h]
	\centering
	
	\tikzstyle{comp} = [circle, draw, minimum size = 0.5cm, line width=0.7] % comparator style
	\tikzstyle{block} = [rectangle, draw, minimum width = 1.2cm, minimum height = 1.2cm, line width=0.7]  % bloc style
	\tikzstyle{big_block} = [rectangle, draw, minimum width = 2cm, minimum height = 2cm, line width=0.7]  % bloc style
	
	\newcommand{\comp}[3][]{ % comparator command
		\node[comp] (#2) at #3 {};
		\draw (#2) -- + (45:0.25) -- + (45:-0.25);
		\draw (#2) -- + (-45:0.25) -- + (-45:-0.25);
		\ifthenelse{
			\isempty{#1}{}
		}{}
		{
			\ifthenelse{
				\equal{#1}{a}
			}{
				\node[above right] at (#2) {$-$\phantom{p}};
			}{}
			\ifthenelse{
				\equal{#1}{b}
			}{
				\node[below right] at (#2) {$-$\phantom{l}};
			}{}
		}
	}
	\subcaptionbox{Schéma de la simulation\label{subfig_impedance_simulation_a}}
	{	
		\tikzsetnextfilename{subfig_impedance_simulation_a}		
		\begin{tikzpicture} [node distance = 2cm]	
			\comp[]{c1}{(0, 0)};
			
			\node[left of = c1, node distance = 1cm] (fz) {};
			
			\node[block, right of = c1] (imp_model_1) {$\frac{1}{K + Bs + Ms^2}$};
			
			\node[block, below of = imp_model_1] (imp_model_2) {$\frac{1}{K + Bs + Ms^2}$};
				
			\draw[-{Latex[length=3mm]}] ($(fz.center) + (-.5,0)$) -- (c1.west)
				node[align = center, above, near start] {$\vect{f}_z$};
			
			\draw[-{Latex[length=3mm]}] (c1.east) -- (imp_model_1.west);
				
			\draw[-{Latex[length=3mm]}] (fz.center) |- (imp_model_2.west);
			
			\draw[-{Latex[length=3mm]}] (imp_model_1.east) --+ (1,0)
				node[align = center, below, midway] {$\vect{x}_z$};
			
			\draw[-{Latex[length=3mm]}] (imp_model_2.east) --+ (1,0)
				node[align = center, below, midway] {$\vect{x}_{z0}$};
			
			\draw[{Latex[length=3mm]}-] (c1.north) --+ (0,.75)
				node[align = center, right] {$\vect{f}_p$};
			
			\fill (fz) circle [radius=2.5pt];
		\end{tikzpicture}
		\vspace{5mm}
	}
	\subcaptionbox{Aspect de la perturbation en force $\vect{f}_p$\label{subfig_impedance_simulation_b}}
	{
		\tikzsetnextfilename{subfig_impedance_simulation_b}		
		\begin{tikzpicture}
			\centering
			%
			\begin{axis}[%
				width=.3\linewidth,
				height=.2\linewidth,
				scale only axis,
				xlabel={temps (\SI{}{\milli\second})},
				ylabel={Force (\SI{}{\newton})},
				axis x line = bottom,
				axis y line = left, 
				ylabel near ticks,
				xmax=150,
				]
				%
				\addplot [
				color=Orient,
				smooth,
				]
				table [col sep=comma, x expr=\coordindex, y=simul_force_pert] {misc/simulated_force_perturbation.csv};
			\end{axis}
		\end{tikzpicture}
	}
	\caption{Simulation pour l'identification des paramètres en impédance avec des signaux en force issus d'expérimentation sans perturbations.}
	\label{fig_impedance_simulation} % Label 
\end{figure}